\scanpage{78}\mspage{38}
\uline{Corollaire} Supposons $\Lambda=\Lambda_0$ i.e. $\Lambda$ est un anneau ``absolu''. Alors tout autom de $P_{M,i}$ qui fixe un repère (i.e. les sommets d'une face) et commute à l'antipodisme est l'identité.

On trouve donc

\uline{Théorème} \unsure{Soient|Soit} $\Lambda\simeq\mathbb Z$ ou $\mathbb Z/n\mathbb Z$, \inserted{$M$ un $\Lambda$-mod. libre de rg 2.} Alors

a) Pour tte structure $\tfrac12$-symplectique $i\in\underbrace{(\det M^*)/\{\pm1\}}_{\unsure{\pi|\Pi}_M}$, \deleted{le groupe des automorphismes} \inserted{l'hom. canonique}
\[
 \mathrm{\deleted{\unsure{S}}Gl}_\Lambda(M)^{\pm}\longrightarrow
 \deleted{\unsure{G_{\widetilde S_M}}}\underbrace{\operatorname{Aut}(\widetilde P_{M,i},\deleted{\unsure{\xi}}\mathrm{antipod})}_{\substack{\text{automs qui}\\\text{commutent à}\\\text{l'antipodisme}}}\tag*{(16)}
\]
\deleted{est injectif (sur unité) et} est un isomorphisme [car $\mathrm{Gl}_\Lambda(M)^{\pm}$ est simplement transitif sur l'ens des bases $\tfrac12$ symplectiques relatives à $i$].

b) \deleted{\unsure{Tout autom de}}\inserted{Considérons sur} $\widetilde S_M=\deleted{\unsure{A}}M^*$ la structure $\widetilde P_M$ définie par la famille $(\widetilde F_{M,i})_{i\in\pi_M}$ \deleted{\unsure{est}} de structures triangulées, et l'antipodisme, qui est un autom \inserted{involutif} commun à toutes ces structures triangulées. Alors
% slanted
\margin{il suffit \unsure{$=$|même} de ne regarder \uline{que} la structure donnée par $\widetilde F_M$ tout seul, \unsure{sans} sa partition}
\[
 \mathrm{Gl}_\Lambda(M)\xrightarrow{\sim}\operatorname{Aut}(\widetilde P_M)\,.\tag*{(17)}
\]
% archive sheet number: 77
